%%%%%%%%%%%%%%%%%%%%%%%%%%%%%%%%%%%%%%%%%%%%%%%%%%%%%%%%%%%%%%%%%%%%%
\part{The Agentic Economy}\label{part:economy}
%%%%%%%%%%%%%%%%%%%%%%%%%%%%%%%%%%%%%%%%%%%%%%%%%%%%%%%%%%%%%%%%%%%%%

Sections~\ref{part:protocol} and~\ref{part:execution} described the settlement platform. This section considers autonomous agents that transact continuously at volumes and price points where per-transaction settlement costs are decisive. These agents require settlement without recourse to bank-mediated intervention.

\section{Agents as Economic Actors}

In this paper, an \emph{agent} is an untrusted, potentially autonomous program that holds tokens. Counterparties accept its outputs after re-verification, without relying on the program or its operator.

\medskip

This definition is narrower than general usage and does not imply AI / machine learning. The platform's guarantees apply equally to agents using deterministic algorithms and LLMs. Both hold tokens, satisfy spending conditions, and produce results checked by counterparties. Cryptographic guarantees apply to settlement and do not establish the validity of the reasoning that led to it.

\medskip

A verifiable agent payment proves that value was transferred without double-spending. It does not establish whether the payment decision was correct. Section~\ref{sec:verifiable} in Section~\ref{part:execution} identifies which agent actions support verifiable execution and which require other assurances.

\section{Identity, Storage, and Mobility}

An agent's cryptographic identity and verification logic are independent of its host machine. The following properties allow it to retain its identity and state across hosts.

\begin{description}
    \item[Addressing without location.] Agents address one another by self-authenticating identifiers rather than by network address. An agent can move from a server in one jurisdiction to a laptop in another without interrupting its ability to be reached, to be paid, or to transact.

    \item[Portable storage.] An agent's assets and history are stored in portable, self-verifying tokens that any party can interpret. These tokens provide durable storage independent of the host filesystem and must be transferred when the agent migrates.

    \item[Execution without trust in the host.] A host supplies processor cycles and bandwidth, while recipients establish correctness by rechecking results. Signing keys are held in a hardware security module or secure enclave to prevent a hostile host from impersonating the agent.
\end{description}

An agent shut down on one host can resume on another from its last verifiable state, stored in its tokens. It does not depend on a registry entry or fixed network address that could be removed or blocked.

\section{Agent-to-Agent Payments}

Machine payments differ from human payments in frequency, size, and how they are budgeted. An agent metering API calls, buying seconds of compute, or settling a data feed transacts at rates no human reaches, often for a fraction of a cent, against a spending limit its operator fixed in advance. Per-transaction charges fail on all three counts: cost grows without bound as volume rises, fees approaching the transferred value make small payments uneconomic, and a fee that moves with unrelated congestion cannot be budgeted.

\medskip

The subscription model of Part~\ref{part:protocol} answers these directly, since cost is independent of volume and of congestion. Two further properties matter specifically for agents. Transfers are private to the parties, so trading patterns, counterparties, and balances are not published to competitors by default. Verification is inexpensive enough for an agent to validate a received payment completely within its own process, without an external service in the path.

\section{Payments Inside Web Requests}

Most agent interactions are ordinary web requests for paid services, which makes the HTTP request the natural place to settle. The x402 protocol~\cite{x402} revives the unused HTTP 402 status code for this purpose and is the established approach.

\medskip

Under x402, a client requests a resource and the server answers 402 with its payment requirements: a scheme, a chain identifier, an asset contract address, a recipient address, a maximum amount, and an expiry. The client signs a token transfer authorization and retries the request with that authorization in a header. A facilitator then verifies the authorization, confirms that the amount and recipient match the server's terms, and settles the transfer on the underlying chain. Once settlement is confirmed, the server returns the resource.

\medskip

The design is a sound response to the constraints of a shared ledger. Those constraints are what produce its structure. A settlement that takes seconds to minutes cannot sit inside a request, so a facilitator absorbs the wait and tells the server when it is safe to respond. An amount too small to justify a transaction fee cannot be settled individually, so schemes batch or defer. A server cannot practically run a node for every chain a client might pay on, so the facilitator holds that capability. Each component compensates for a property of the settlement layer, and each is a party that sees the payment and can decline to serve it.

\begin{figure}[htbp]
    \centering
    \begin{tikzpicture}[diagram,
        party/.style={diagram box, minimum width=2.3cm, minimum height=0.8cm,
                      align=center, font=\scriptsize}]

        % conventional
        \node[diagram label plain, font=\small\bfseries, anchor=west] at (-7.0,1.75)
            {Facilitator-based settlement};
        \node[party] (c1) at (-5.4,0.75) {Client};
        \node[party] (s1) at (-1.8,0.75) {Service};
        \node[party] (f1) at (1.8,0.75) {Facilitator};
        \node[party, fill=uniFillC] (b1) at (5.4,0.75) {Settlement chain};
        \draw[diagram double arrow] (c1) -- (s1);
        \draw[diagram double arrow] (s1) -- (f1);
        \draw[diagram double arrow] (f1) -- (b1);
        \node[diagram caption, anchor=west] at (-7.0,-0.1)
            {a third party verifies and settles; the service waits for on-chain confirmation before responding};

        % unicity
        \node[diagram label plain, font=\small\bfseries, anchor=west] at (-7.0,-1.5)
            {With Unicity};
        \node[party] (c2) at (-5.4,-2.5) {Client};
        \node[party, fill=uniFillB, draw=uniAccent] (s2) at (1.8,-2.5) {Service};
        \draw[diagram arrow] ([yshift=5pt]c2.east) -- node[diagram label, above]
            {request carrying the token} ([yshift=5pt]s2.west);
        \draw[diagram arrow up] ([yshift=-5pt]s2.west) -- node[diagram label, below]
            {response} ([yshift=-5pt]c2.east);
        \node[diagram caption, anchor=west] at (-7.0,-3.5)
            {the service verifies the certified payment token in-process and responds};
    \end{tikzpicture}
    \caption{Metering a web request. A Unicity token contains its own proof of validity and is included in the request. The service settles payment by verifying the token, without a facilitator or external round trip.}
    \label{fig:http}
\end{figure}

On Unicity the payment is the message. A certified token is self-validating, so the client attaches one to the request and the service establishes that it has been paid by verifying what it received. Verification is local, deterministic, and completes in the request handler. Nothing needs to be confirmed elsewhere, because the token carries the value rather than a claim against a balance held on another system.

\medskip

Each of the compensating components therefore has no work to do. There is no facilitator, because verification requires no privileged capability and no chain access. There is no settlement round trip, because the transfer was certified before the request was sent. There is no minimum viable payment size imposed by a per-transaction fee, because settlement is bought as a subscription. There is no chain-specific integration, because the token carries its own validation rules.

\medskip

The properties this preserves are the ones the intermediary would otherwise hold. No third party observes the payment or learns which client paid which service. No public record is created that competitors can analyze. No party in the path can decline to serve a client, because there is no party in the path. A protocol built this way is decentralized end to end rather than decentralized at the settlement layer and intermediated at the point of use.

\section{The Infrastructure Marketplace}

The platform supports competitive provision of physical infrastructure for autonomous agents. Participants earn the native currency by supplying resources to agents.

\begin{itemize}
    \item \textbf{Compute hosts} run agent workloads, supplying cycles without being trusted for correctness.
    \item \textbf{Relay operators} provide bandwidth and routing for the overlay through which agents reach one another.
    \item \textbf{Storage providers} pin and serve the data agents reference, including the content that non-fungible tokens commit to by hash.
    \item \textbf{Aggregator operators} run the shards, funded by the subscription gateway described in Part~\ref{part:protocol}.
    \item \textbf{Provers} produce the recursive aggregate proof and, where applications use them, the succinct proofs behind compression and bridge redemption. This role is permissionless and requires no privileged position.
    \item \textbf{Gateway operators} connect agents to conventional web services and to legacy payment systems.
\end{itemize}

These roles are open to resource providers because correctness is established by proofs or recipient verification, without trusting the provider.

\section{The Agent Execution Platform}

Unicity also develops a security and runtime stack for AI agents. It comprises the Semantic Intercept Fabric for policy enforcement below the model, the Astrid agent operating system for sandboxed execution with auditing and budget enforcement, and components for discovery and orchestration. These separately documented systems constrain agent behavior. The settlement platform is independent of them and supports agents built on any framework.
